\chapter{Résultats}

\section{Fonctionnalités réalisées}
Au moment de la rédaction, Secure Requirements Studio permet~:
\begin{itemize}
    \item de mener un entretien de cadrage complet, en mode Classic Questionnaire ou
    assisté par IA (Gemini) ;
    \item de générer dynamiquement des exigences fonctionnelles, non fonctionnelles, de
    sécurité et techniques, traçables aux réponses de l'entretien ;
    \item de produire une analyse de sécurité (score, checklist, matrice RBAC, registre
    des risques, analyse d'impact vie privée) ;
    \item d'exécuter une revue de complétude avant export ;
    \item de générer un Cahier des Charges et une conception applicative avec un
    diagramme MVP unique, exportables en JSON, Markdown, HTML, DOCX et PDF.
\end{itemize}

\section{Résultats obtenus}
\TODO{Compléter avec des résultats concrets une fois un ou plusieurs entretiens réels
menés avec des parties prenantes de la Fondation OCP — ne pas anticiper de résultats
métier qui n'ont pas encore été produits.}

\section{Limites}
\begin{itemize}
    \item SRS produit une \emph{conception et une spécification}, pas une implémentation
    de la future Plateforme de Gestion des Bénéficiaires elle-même.
    \item L'intégration de certains fournisseurs IA envisagés initialement (OpenAI,
    Anthropic Claude, Kimi, Ollama) peut être partielle ou absente selon l'état réel du
    code au moment de la rédaction — voir chapitre~6.
    \item L'outil ne dispose volontairement d'aucune authentification, ce qui limite
    son usage à un contexte interne restreint.
    \item \TODO{Ajouter toute autre limite constatée en conditions réelles d'usage.}
\end{itemize}

\section{Perspectives}
\begin{itemize}
    \item Étendre la couverture des fournisseurs IA supportés.
    \item Enrichir le catalogue de concepts du Knowledge Graph au fur et à mesure des
    entretiens réels menés avec les parties prenantes de l'Axe Économie Sociale et
    Solidaire.
    \item Transmettre les livrables générés (Cahier des Charges, conception, diagramme
    MVP, exigences de sécurité) à l'équipe de développement chargée de la future
    Plateforme de Gestion des Bénéficiaires.
\end{itemize}
